\uuid{VlZr}
\titre{Fonction majorée/minorée}
\auteur{bodin ; hanani ; mzari}
\organisation{exo7}
\licence{CC BY-NC-SA 4.0}
\datecreate{2020-01-20}
\niveau{L1}
\module{Analyse}
\chapitre{Continuité, limite et étude de fonctions réelles}
\sousChapitre{Autre}
\difficulte{2}
\theme{Notion de fonctions, fonction majorée, minorée}
\exo7qcmid{lille-1.continuite.005}

\contenu{
\begin{question}
\qclassification{123.00}
\qtags{motcle=fonction majorée/minorée}
Soit $f : \Rr \to \Rr$. Quelles sont les assertions vraies ?
\begin{answers}
    \bad{$\forall  M > 0 \quad \exists x \in \Rr \quad f(x) \le M$  implique $f$ majorée.}

    \bad{$\forall  x \in \Rr \quad \exists  M > 0 \quad f(x) \ge M$  implique $f$ majorée.}

    \good{$\exists  M > 0 \quad \forall x \in \Rr \quad f(x) \le M$  implique $f$ majorée.}

    \bad{$\exists  M > 0 \quad \exists x \in \Rr \quad f(x) \ge M$  implique $f$ majorée.}

\end{answers}
\begin{explanations}
Par définition $f$ est majorée si $\exists  M > 0 \quad \forall x \in \Rr \quad f(x) \le M$.
\end{explanations}
\end{question}
}
